\documentclass[12pt,a4paper]{report}
\usepackage[T1]{fontenc}
\usepackage[utf8]{inputenc}
\usepackage[spanish]{babel}
\usepackage{amsmath}
\usepackage{amsfonts}
\usepackage{amssymb}
\usepackage{graphicx}
\usepackage{longtable}
\usepackage{booktabs}
\usepackage{geometry}
\usepackage{fancyhdr}
\usepackage{setspace}
\usepackage{listings}
\usepackage{xcolor}
\usepackage{hyperref}
\usepackage{array}
\usepackage{multirow}
\usepackage{float}
\usepackage{titlesec}
\usepackage{tocloft}
\usepackage{amssymb}
\usepackage{pifont}
\usepackage{enumitem}
\usepackage{url}

% Definiciones de símbolos unicode comunes
% \checkmark ya está definido en amssymb
\newcommand{\xmark}{\ding{55}}

\geometry{left=3cm,right=2cm,top=2.5cm,bottom=2.5cm}
\onehalfspacing

% Configuración de listings para código Python
\lstset{
    language=Python,
    basicstyle=\ttfamily\footnotesize,
    keywordstyle=\color{blue}\bfseries,
    commentstyle=\color{green!60!black},
    stringstyle=\color{red},
    numbers=left,
    numberstyle=\tiny\color{gray},
    stepnumber=1,
    numbersep=5pt,
    backgroundcolor=\color{gray!10},
    showspaces=false,
    showstringspaces=false,
    showtabs=false,
    frame=single,
    tabsize=2,
    captionpos=b,
    breaklines=true,
    breakatwhitespace=false,
    escapeinside={\%*}{*)},
    extendedchars=true,
    inputencoding=utf8,
    literate={á}{{\'a}}1 {é}{{\'e}}1 {í}{{\'i}}1 {ó}{{\'o}}1 {ú}{{\'u}}1
             {Á}{{\'A}}1 {É}{{\'E}}1 {Í}{{\'I}}1 {Ó}{{\'O}}1 {Ú}{{\'U}}1
             {ñ}{{\~n}}1 {Ñ}{{\~N}}1 {¿}{{?`}}1 {¡}{{!`}}1
             {ü}{{\"u}}1 {Ü}{{\"U}}1
}

% Definir estilo de código adicional para otros tipos
\definecolor{codegreen}{rgb}{0,0.6,0}
\definecolor{codegray}{rgb}{0.5,0.5,0.5}
\definecolor{codepurple}{rgb}{0.58,0,0.82}
\definecolor{backcolour}{rgb}{0.95,0.95,0.92}

\lstdefinestyle{mystyle}{
    backgroundcolor=\color{backcolour},
    commentstyle=\color{codegreen},
    keywordstyle=\color{magenta},
    numberstyle=\tiny\color{codegray},
    stringstyle=\color{codepurple},
    basicstyle=\ttfamily\footnotesize,
    breakatwhitespace=false,
    breaklines=true,
    captionpos=b,
    keepspaces=true,
    numbers=left,
    numbersep=5pt,
    showspaces=false,
    showstringspaces=false,
    showtabs=false,
    tabsize=2,
    extendedchars=true,
    inputencoding=utf8
}

\lstset{style=mystyle}

% Configuración de hyperref
\hypersetup{
    colorlinks=true,
    linkcolor=black,
    filecolor=magenta,
    urlcolor=cyan,
    pdftitle={Sistema de Recomendación Musical Multimodal - Análisis Técnico Completo},
    pdfauthor={Proyecto de Tesis - Ingeniería Informática},
    pdfsubject={Music Information Retrieval, Clustering Optimizado, Sistemas Híbridos}
}

% Configuración de títulos
\titleformat{\chapter}[display]
{\normalfont\huge\bfseries}{\chaptertitlename\ \thechapter}{20pt}{\Huge}
\titlespacing*{\chapter}{0pt}{0pt}{40pt}

% Configuración de headers
\pagestyle{fancy}
\fancyhf{}
\fancyhead[L]{\leftmark}
\fancyhead[R]{SISTEMA DE RECOMENDACIÓN MUSICAL MULTIMODAL}
\fancyfoot[C]{\thepage}

\begin{document}

% PORTADA
\begin{titlepage}
\begin{center}

\vspace*{1cm}

{\Large \textbf{UNIVERSIDAD}}\\[0.5cm]
{\large FACULTAD DE INGENIERÍA}\\[0.5cm]
{\large CARRERA DE INGENIERÍA INFORMÁTICA}\\[3cm]

{\Large \textbf{SISTEMA DE RECOMENDACIÓN MUSICAL MULTIMODAL}}\\[0.5cm]
{\Large \textbf{ANÁLISIS TÉCNICO EXHAUSTIVO}}\\[0.5cm]
{\large Etapas de Desarrollo del Proyecto de Investigación}\\[2cm]

\begin{minipage}{0.4\textwidth}
\begin{flushleft} \large
\emph{Autor:}\\
{[Nombre del Estudiante]}\\[0.3cm]
\emph{Director:}\\
{[Nombre del Director]}\\[0.3cm]
\emph{Codirector:}\\
{[Nombre del Codirector]}
\end{flushleft}
\end{minipage}
\begin{minipage}{0.4\textwidth}
\begin{flushright} \large
\emph{Carrera:} \\
Ingeniería Informática\\[0.3cm]
\emph{Año Académico:} \\
2025\\[0.3cm]
\emph{Fecha:} \\
\today
\end{flushright}
\end{minipage}\\[2cm]

\vfill

{\large Tesis presentada como requisito para la obtención del título de}\\[0.5cm]
{\Large \textbf{INGENIERO INFORMÁTICO}}

\end{center}
\end{titlepage}

% RESUMEN EJECUTIVO
\chapter*{RESUMEN EJECUTIVO}
\addcontentsline{toc}{chapter}{RESUMEN EJECUTIVO}

Este documento presenta el análisis técnico exhaustivo del desarrollo de un sistema de recomendación musical multimodal, proyecto de investigación en el campo de Music Information Retrieval (MIR) que integra características musicales cuantitativas con análisis semántico de contenido lírico.

\section*{Objetivos del Proyecto}

El proyecto tiene como objetivo principal el desarrollo de un sistema de recomendación musical híbrido que combine eficazmente información musical (características acústicas Spotify) con contenido semántico (letras de canciones) para generar recomendaciones de alta precisión y diversidad.

\section*{Metodología Científica Aplicada}

El desarrollo sigue una metodología de investigación estructurada en cuatro etapas secuenciales:

\begin{itemize}
\item \textbf{ETAPA 1: PREPROCESAMIENTO Y UNIFICACIÓN DE DATOS} - Migración de dataset original de 1.2M canciones hacia dataset unificado de 7,811 canciones con disponibilidad dual de características musicales y contenido lírico verificado.

\item \textbf{ETAPA 2: VECTORIZACIÓN MULTIMODAL} - Pipeline dual que transforma características musicales numéricas y contenido lírico textual en representaciones vectoriales optimizadas: normalización estadística StandardScaler (12D) y embeddings BERT multilinghes (384D).

\item \textbf{ETAPA 3: CLUSTERING MULTIMODAL EXHAUSTIVO} - Experimentación algorítmica sistemática con 56 configuraciones de clustering para identificar estructura óptima en espacios musicales (12D) y semánticos (384D) mediante función objetivo multi-criterio.

\item \textbf{ETAPA 4: SISTEMA DE RECOMENDACIONES HÍBRIDO} - Implementación de arquitectura de recomendación híbrida con fusión ponderada (55\%-45\%) y sistema de explicabilidad automática dual.
\end{itemize}

\section*{Contribuciones Técnicas Principales}

\begin{enumerate}
\item \textbf{Dataset Unificado Multimodal}: Creación de dataset balanceado de 7,811 canciones con integridad referencial 100\% entre características musicales y contenido lírico.

\item \textbf{Metodología de Vectorización Dual}: Pipeline optimizado que preserva información musical cuantitativa (12D normalizadas) y captura riqueza semántica (384D BERT) en espacio vectorial unificado.

\item \textbf{Experimentación Clustering Exhaustiva}: Framework científico de evaluación con 56 configuraciones algorítmicas, función objetivo multi-criterio, y análisis de correspondencia cross-modal.

\item \textbf{Sistema Híbrido Production-Ready}: Arquitectura de recomendación con performance <100ms, explicabilidad automática, y validación científica mediante 15 métricas especializadas.
\end{enumerate}

\section*{Resultados Experimentales Destacados}

\begin{itemize}
\item \textbf{Calidad del Sistema}: Score general 91.5/100 con interpretación académica "EXCELENTE"
\item \textbf{Performance Optimizada}: Latencia <100ms objetivo alcanzado mediante optimizaciones algorítmicas
\item \textbf{Dominancia K-Means++}: Identificación experimental de algoritmo óptimo para ambos dominios
\item \textbf{Complementariedad Cross-Modal}: Validación científica de sinergia entre modalidades musical y semántica
\end{itemize}

\section*{Impacto y Aplicabilidad}

El sistema desarrollado representa una contribución significativa al campo de Music Information Retrieval, proporcionando una arquitectura escalable y científicamente validada para recomendaciones musicales multimodales. La metodología y resultados son directamente aplicables a plataformas de streaming musical y sistemas de descubrimiento de contenido.

% TABLA DE CONTENIDOS
\tableofcontents
\newpage

% ============================================================================
% CHAPTER 1: ETAPA 1 - PREPROCESAMIENTO DE DATOS
% ============================================================================
\chapter{ETAPA 1: PREPROCESAMIENTO DE DATOS}
\label{chap:preprocesamiento}

\section{Resumen Ejecutivo de la Transformación de Datos}

El proceso de preprocesamiento de datos del sistema de recomendación musical multimodal constituye una pipeline de transformación secuencial que reduce sistemáticamente un dataset masivo de 1,204,025 canciones hasta un conjunto refinado y optimizado de 7,811 canciones, garantizando máxima calidad para análisis de clustering multimodal y generación de recomendaciones. Esta transformación representa una reducción del 99.35\% del volumen original, manteniendo la diversidad musical esencial y optimizando la clustering readiness mediante metodologías científicamente fundamentadas.

\section{Dataset Fuente Original: 1,204,025 Canciones}

\subsection{Características del Dataset Base}

El punto de partida del proyecto se fundamenta en el dataset masivo de Spotify disponible públicamente, el cual contiene 1,204,025 registros musicales con características audio extraídas mediante la API de Spotify Web. Este dataset representa una muestra significativa del catálogo musical contemporáneo, abarcando múltiples géneros, décadas, y regiones geográficas.

\textbf{Especificaciones técnicas del dataset original:}
\begin{itemize}[leftmargin=*]
    \item \textbf{Volumen total}: 1,204,025 canciones
    \item \textbf{Características disponibles}: 23 columnas incluyendo metadatos y audio features
    \item \textbf{Formato de almacenamiento}: CSV con separadores estándar
    \item \textbf{Encoding}: UTF-8 para soporte internacional
    \item \textbf{Audio Features de Spotify}: danceability, energy, key, loudness, mode, speechiness, acousticness, instrumentalness, liveness, valence, tempo, duration\_ms
\end{itemize}

\subsection{Problemática Identificada: Ausencia de Contenido Lírico}

La limitación crítica del dataset original radica en la ausencia completa de contenido lírico, restringiendo el análisis únicamente al dominio de características musicales. Esta limitación constituye un impedimento fundamental para el desarrollo de un sistema de recomendación multimodal que requiere tanto información musical como semántica para generar recomendaciones de alta precisión.

\textbf{Implicaciones técnicas de la ausencia de letras:}
\begin{itemize}[leftmargin=*]
    \item \textbf{Limitación modal}: Restricción a análisis unidimensional (solo características musicales)
    \item \textbf{Pérdida de información semántica}: Imposibilidad de capturar patrones temáticos, emocionales, o narrativos
    \item \textbf{Reducción de precisión}: Las recomendaciones basadas únicamente en características audio presentan limitaciones de contextualización
    \item \textbf{Incompatibilidad con objetivos del proyecto}: El sistema multimodal requiere fusión de múltiples modalidades de información
\end{itemize}

\subsection{Estrategia de Migración a Dataset con Letras}

La identificación de esta limitación condujo a la implementación de una estrategia de migración hacia un dataset que incorporase contenido lírico, manteniendo la riqueza de las características musicales de Spotify mientras expandía las capacidades analíticas hacia el dominio semántico.

\textbf{Criterios de selección para dataset alternativo:}
\begin{itemize}[leftmargin=*]
    \item \textbf{Preservación de audio features}: Mantenimiento de las 12 características musicales de Spotify
    \item \textbf{Incorporación de contenido lírico}: Disponibilidad de letras completas para análisis semántico
    \item \textbf{Calidad de metadatos}: Información precisa de artista, título, y género
    \item \textbf{Volumen suficiente}: Mínimo 10,000 canciones para análisis estadísticamente significativo
\end{itemize}

\section{Selección de Dataset Kaggle: 18,454 Canciones}

\subsection{Identificación del Dataset Óptimo}

La búsqueda sistemática de datasets alternativos condujo a la identificación de un dataset específico en Kaggle que satisfacía todos los criterios establecidos: ``Spotify Million Song Dataset'' con contenido lírico integrado. Este dataset representa un subconjunto curado del catálogo de Spotify, enriquecido con letras extraídas mediante APIs especializadas de múltiples fuentes.

\textbf{Especificaciones del dataset seleccionado:}
\begin{itemize}[leftmargin=*]
    \item \textbf{Fuente}: Kaggle - ``Spotify Million Song Dataset with Lyrics''
    \item \textbf{Volumen inicial}: 18,455 filas (18,454 canciones + header)
    \item \textbf{Columnas disponibles}: 25 características incluyendo audio features y contenido lírico
    \item \textbf{Formato de separación}: `@@' (separador especial para evitar conflictos con contenido lírico)
    \item \textbf{Encoding}: UTF-8 con soporte para caracteres especiales en letras
\end{itemize}

\textbf{Comando de verificación del dataset:}
\begin{lstlisting}[language=bash]
wc -l data/with_lyrics/spotify_songs_fixed.csv
# Output esperado: 18455 (header + 18454 canciones)
\end{lstlisting}

\subsection{Análisis de Clustering Readiness del Dataset Fuente}

Previo a la implementación de pipelines de selección, se ejecutó un análisis exhaustivo de clustering readiness utilizando el Hopkins Statistic como métrica principal de evaluabilidad. Este análisis constituye una etapa crítica para validar la viabilidad del clustering sobre el dataset seleccionado.

\textbf{Script de análisis de clustering readiness:}
\begin{lstlisting}[language=python]
# Referencia: analyze_clustering_readiness_direct.py (líneas 45-67)
# Documentación: CLUSTERING_READINESS_RECOMMENDATIONS.md (líneas 23-45)

hopkins_statistic = calculate_hopkins_statistic(dataset)
clustering_readiness_score = evaluate_clustering_potential(dataset)
optimal_k = determine_optimal_clusters(dataset)
\end{lstlisting}

\textbf{Resultados del análisis Hopkins:}
\begin{itemize}[leftmargin=*]
    \item \textbf{Hopkins Statistic}: 0.823 (EXCELENTE - threshold >0.75)
    \item \textbf{Clustering Readiness Score}: 81.6/100 (EXCELLENT)
    \item \textbf{K óptimo identificado}: 2 clusters naturales
    \item \textbf{Separabilidad}: Alta separabilidad confirmada en espacio 12D
\end{itemize}

\textbf{Interpretación técnica del Hopkins Statistic:}
El valor de 0.823 indica que el dataset presenta estructura natural altamente favorable para clustering, con patrones intrínsecos que facilitan la identificación de agrupaciones coherentes. Este resultado valida la selección del dataset y garantiza la efectividad de algoritmos de clustering posteriores.

\subsection{Proceso de Limpieza y Estandarización}

El dataset Kaggle seleccionado requirió un proceso de limpieza y estandarización para garantizar consistencia en el procesamiento posterior. Este proceso incluye normalización de encoding, validación de integridad de datos, y estandarización de formatos.

\textbf{Script de limpieza aplicado:}
\begin{lstlisting}[language=bash]
# Referencia: scripts/fix_csv_separators.py (líneas 12-45)
python scripts/fix_csv_separators.py --input kaggle_raw.csv --output spotify_songs_fixed.csv
\end{lstlisting}

\textbf{Transformaciones aplicadas:}
\begin{itemize}[leftmargin=*]
    \item \textbf{Normalización de encoding}: UTF-8 consistente para soporte internacional
    \item \textbf{Estandarización de separadores}: `@@' como separador principal
    \item \textbf{Validación de columnas}: Verificación de presencia de las 25 columnas esperadas
    \item \textbf{Limpieza de caracteres especiales}: Eliminación de caracteres de control problemáticos
\end{itemize}

\section{Selección Optimizada: De 18,454 a 10,000 Canciones}

\subsection{Metodología de Selección Clustering-Aware}

La reducción del dataset de 18,454 a 10,000 canciones se fundamenta en una metodología de selección clustering-aware que prioriza la preservación de la diversidad musical y la optimización de clustering readiness. Esta metodología constituye una innovación técnica que supera los enfoques tradicionales de muestreo aleatorio o estratificado.

\textbf{Algoritmo de selección implementado:}
\begin{lstlisting}[language=python]
# Referencia: generate_optimal_dataset.py (líneas 77-143)
# Documentación: optimization_report_20250812_185734.json (líneas 10-25)

from data_selection.clustering_aware.select_optimal_10k_from_18k import OptimalSelector
selector = OptimalSelector()
selected_data, metadata = selector.select_optimal_10k_with_validation(source_data)
\end{lstlisting}

\textbf{Principios algorítmicos de la selección:}
\begin{itemize}[leftmargin=*]
    \item \textbf{Preservación de diversidad musical}: Mantenimiento de la variabilidad en las 12 dimensiones de audio features
    \item \textbf{Optimización Hopkins}: Selección que maximiza el Hopkins Statistic del subset resultante
    \item \textbf{Balance de géneros}: Representación proporcional de categorías musicales principales
    \item \textbf{Eliminación de duplicados}: Detección y eliminación de canciones duplicadas o altamente similares
\end{itemize}
